\clearpage\fancyhead[L]{\small GAUSS AN BODE.}\hypertarget{native-404-part-lower}{}\label{native-404-part-lower}
\begin{center}\large GAUSS AN BODE.\end{center}
\begin{center}\hrule\vspace{4pt}\fontsize{10}{12}\selectfont
Astronomisches Jahrbuch für 1820. Seite 201..205. Berlin 1817.\par\vspace{4pt}\hrule\vspace{4pt}\end{center}
\begin{flushright}Göttingen, den 17. März 1817.\end{flushright}
\par\indent Da die Theorie der Pallasbewegung hinlänglich ausgebildet ist, so habe ich Hrn. Dr. \textsc{Tittel} aus Erlau, welcher sich anderthalb Jahre hier aufgehalten, veranlasst, die Ephemeride für 1818..1819 schon \SourceLeaf{405} jetzt zu berechnen. Ich habe das Vergnügen, Ihnen dieselbe hier für Ihr nächstes Jahrbuch in Abschrift beizulegen.

\par\indent Ist der Himmel im vorigen Winter bei Ihnen den Beobachtungen eben so ungünstig gewesen, wie hier? Höchst selten war die Sonne oder ein Stern zu erblicken. Am ganzen 19. Nov. v.\,J. schneiete es, und von der Sonnenfinsterniss war nichts zu sehen. Einige Abnahme der Tagshelle war um die Mitte der Finsterniss zu spüren, die jedoch nichts ausserordentliches hatte; wenigstens da es seitdem, selbst um Mittag, bei Schneegestöber, auch ohne Sonnenfinsterniss, öfters eben so dunkel gewesen. Bei der Mondfinsterniss vom 4. Dec. war der Himmel etwas weniger ungünstig. Anfang und Ende blieben zwar unsichtbar, allein folgende Phasen habe ich zum Theil zwischen fliegenden Wolken ganz gut beobachtet.

\Needspace{15\baselineskip}\par\indent 1816. December 4 auf der neuen Sternwarte.

\par\indent Eintritte.
\begin{center}\fontsize{10}{12}\selectfont\setlength{\tabcolsep}{3pt}\begin{tabular}{llrrrl}
Tycho 1. Rand & $\cdot\;\cdot\;\cdot\;\cdot\;\cdot\;\cdot\;\cdot\;\cdot\;\cdot\;\cdot\;\cdot\;\cdot$ &  $8^{\mathrm u}$ & $17^{\mathrm m}$ & $2^{\mathrm s}$ & W. Z.\\
Tycho 2. Rand & $\cdot\;\cdot\;\cdot\;\cdot\;\cdot\;\cdot\;\cdot\;\cdot\;\cdot\;\cdot\;\cdot\;\cdot$ &  & $18$ & $44$ & —\\
Kepler Mitte & $\cdot\;\cdot\;\cdot\;\cdot\;\cdot\;\cdot\;\cdot\;\cdot\;\cdot\;\cdot\;\cdot\;\cdot\;\cdot$ &  & $29$ & $24$ & —\\
Copernicus 1. Rand & $\cdot\;\cdot\;\cdot\;\cdot\;\cdot\;\cdot\;\cdot\;\cdot\;\cdot\;\cdot\;\cdot$ &  & $40$ & $28$ & —\\
Copernicus 2. Rand & $\cdot\;\cdot\;\cdot\;\cdot\;\cdot\;\cdot\;\cdot\;\cdot\;\cdot\;\cdot\;\cdot$ &  & $45$ & $11$ & —\\
Aristarch 1. Rand & $\cdot\;\cdot\;\cdot\;\cdot\;\cdot\;\cdot\;\cdot\;\cdot\;\cdot\;\cdot\;\cdot$ &  & $50$ & $28$ & —\\
Dionysius Mitte & $\cdot\;\cdot\;\cdot\;\cdot\;\cdot\;\cdot\;\cdot\;\cdot\;\cdot\;\cdot\;\cdot\;\cdot$ &  $9$ & $0$ & $46$ & —\\
Censorinus Mitte & $\cdot\;\cdot\;\cdot\;\cdot\;\cdot\;\cdot\;\cdot\;\cdot\;\cdot\;\cdot\;\cdot$ &  & $5$ & $44$ & —\\
\end{tabular}\end{center}

\Needspace{7\baselineskip}\begin{flushright}Göttingen, den 15. August 1817.\end{flushright}

\par\indent Der Bau der hiesigen neuen Sternwarte ist zwar noch nicht ganz, aber doch so weit vollendet, dass in zwei Zimmern beobachtet werden kann. Bis zur Ankunft und Aufstellung der neuen festen Meridianinstrumente werden die Beobachtungen mit dem 12z. \textsc{Reichenbach}schen Wiederholungskreise das Erheblichste sein. Ich theile Ihnen davon die Resultate der Beobachtungen des Nordsterns, und der Sonne im letzten Solstitium mit. Zuvörderst muss ich bemerken, dass nach meinen geodätischen Messungen die Mitte der neuen Sternwarte $6''1$ südlicher, und $1^{\mathrm s}90$ in Zeit östlicher liegt, als die Mitte der alten Sternwarte. Ferner, dass ich an dem Instrument selbst eine kleine Abänderung vorgenommen habe. Bei meinen frühern Beobachtungen bediente ich mich durchgehends des messingenen Ringes, der an der Objectivseite des Fernrohrs vorgesteckt wird, und nicht blos das Uebergewicht des prismatischen Oculars, sondern zugleich die Klammer balancirt. Man hatte vermuthet, dass durch diesen Ring eine Biegung des Fernrohrs (oder genauer zu reden, eine grössere Biegung der Objectivseite desselben, als durch das Ocular an der andern Seite erfolgt) hervorgebracht werden könnte. Die Folge davon würde sein, dass das Instrument alle Zenithdistanzen zu \emph{klein} geben müsste, und man muthmasste, dass dies vielleicht die Ursache von dem Unterschiede sei, den man bei Bestimmung der Polhöhe findet, je nachdem man Circumpolarsterne oder die Sonne beobachtet. Wäre jene Vermuthung gegründet, so müsste der Unterschied verschwinden, wenn man jenes Gegengewicht wegnimmt, oder vielmehr der Unterschied müsste dadurch ins Entgegengesetzte übergehen, da entschieden jetzt die Ocularseite des Fernrohrs ein Uebergewicht hat. Allein meine neuen Beobachtungen, wo jenes Gewicht weggenommen, und dagegen ein Gewicht an den Alhidadenkreis selbst geschraubt ist, sprechen ganz dagegen, und scheinen mir vollkommen zu beweisen, dass die Biegung des Fernrohrs durchaus unmerklich ist. Theils geben die neuen Beobachtungen des Nordsterns und der Sonne wieder ganz denselben Unterschied der Pol\SourceLeaf{406}höhe, und in demselben Sinn, wie die andern; theils ist die absolute Bestimmung der Polhöhe der
neuen Sternwarte ganz einerlei mit der von der alten Sternwarte übertragenen.

\Needspace{15\baselineskip}\par\indent\emph{Beobachtungen des Nordsterns in der untern Culmination.}

\begin{center}\fontsize{10}{12}\selectfont\setlength{\tabcolsep}{3pt}
\begin{tabular}{l|c|c|c}
1817 & Anzahl & \shortstack{Wahre Zenithdi-\\stanz im Meridian} & Polhöhe \\
\hline
März 28 & 16 & $40^\circ\ 8'\ 13''18$ & $51^\circ\ 31'\ 49''33$ \\
April\ \ 1 & 22 & $15.86$ & $47.84$ \\
\ \ \ \ \ \ \ \ 2 & 24 & $13.87$ & $50.13$ \\
\ \ \ \ \ \ \ \ 8 & 28 & $15.76$ & $50.03$ \\
\ \ \ \ \ \ \ 11 & 26 & $16.16$ & $50.51$ \\
\ \ \ \ \ \ \ 13 & 12 & $20.51$ & $46.75$ \\
Mai\ \ \ 7 & 30 & $23.90$ & $49.60$ \\
Mittel & 158 & $\cdot\ \cdot\ \cdot\ \cdot\ \cdot$ & $51^\circ\ 31'\ 49''42$ \\
\end{tabular}
\end{center}

\par\indent Bei Berechnung der Refraction ist \textsc{Carlini}'s Tafel, für die Declination des Nordsterns die
\textsc{Bessel}sche zum Grunde gelegt. Der Beobachtungsplatz liegt $0''12$ nördlich von der Mitte; die
Polhöhe von dieser ist also, aus obigen Beobachtungen $51^\circ\ 31'\ 49''3$.

\par\indent Meine Beobachtungen des Nordsterns in der untern Culmination von 1813 geben mit denselben
Elementen reducirt die Polhöhe der alten Sternwarte $51^\circ\ 31'\ 55''48$; die Uebertragung
gibt also dasselbe was die unmittelbaren Beobachtungen geben.

\par\indent Die Beobachtungen des Sommer-Sonnenstillstandes wurden in diesem Jahre von der Witterung mehr
als je begünstigt: der Beobachtungsplatz liegt $0''16$ südlich von der Mitte.

\Needspace{23\baselineskip}\par\indent\emph{Beobachtungen des Sonnenstillstandes.}

\begin{center}\fontsize{10}{12}\selectfont\setlength{\tabcolsep}{3pt}
\begin{tabular}{l|c|c|c}
1817 & Anzahl & \shortstack{Wahre Zenithdi-\\stanz im Meridian} & Im Solstitium \\
\hline
Juni\ 13 & 10 & $28^\circ\ 18'\ 11''30$ & $28^\circ\ 3'\ 52''26$ \\
\ \ \ \ \ \ 16 & 16 & $9\ \ 51.80$ & $48.63$ \\
\ \ \ \ \ \ 17 & 20 & $7\ \ 57.26$ & $55.30$ \\
\ \ \ \ \ \ 18 & 16 & $6\ \ 19.03$ & $53.46$ \\
\ \ \ \ \ \ 19 & 12 & $5\ \ \ 8.69$ & $54.66$ \\
\ \ \ \ \ \ 20 & 20 & $4\ \ 18.74$ & $51.41$ \\
\ \ \ \ \ \ 21 & 12 & $3\ \ 54.41$ & $48.91$ \\
\ \ \ \ \ \ 22 & 10 & $4\ \ \ 0.64$ & $52.10$ \\
\ \ \ \ \ \ 23 & 12 & $4\ \ 28.17$ & $51.77$ \\
\ \ \ \ \ \ 24 & 16 & $5\ \ 20.62$ & $51.62$ \\
\ \ \ \ \ \ 25 & 10 & $6\ \ 36.39$ & $49.96$ \\
\ \ \ \ \ \ 26 & 14 & $8\ \ 20.16$ & $51.61$ \\
\ \ \ \ \ \ 27 & 22 & $10\ 28.75$ & $53.49$ \\
\ \ \ \ \ \ 29 & 16 & $15\ 54.34$ & $51.86$ \\
\ \ \ \ \ \ 30 & 8 & $19\ 16.59$ & $53.88$ \\
\multicolumn{3}{l}{Mittel aus 214 Beobachtungen} & $28^\circ\ 3'\ 52''15$ \\
\end{tabular}
\end{center}

\par\indent Die Zenithdistanz in der ersten Columne ist blos auf den Meridian reducirt und von der
Refraction befreit; die letzte Columne ist von Parallaxe und dem Einfluss der Sonnenbreite
befreit und auf den Augenblick der Sonnenwende reducirt. Für den Mittelpunkt der Sternwarte
erhalte ich also $28^\circ\ 3'\ 52''31$. Mit der Schiefe der Ekliptik aus den Mailänder
Ephemeriden $23^\circ\ 27'\ 52''9$ folgt hieraus die Polhöhe $51^\circ\ 31'\ 45''21$, um $4''1$
kleiner, als aus dem Polarstern, oder umgekehrt, wenn die Polhöhe aus dem Polarstern zum Grunde
gelegt wird, folgt die Schiefe $4''1$ grösser als nach \textsc{Carlini}, nach \textsc{Bessel}
würde der Unterschied noch anderthalb Secunden grösser sein. Nach Allem, was bisher über dieses \SourceLeaf{407} merkwürdige Phänomen gesagt ist, scheint es nicht, dass wir der Erklärung
desselben näher gekommen wären. Wünschenswerth wäre es gewiss, wenn einmal
unter einer nicht zu geringen Breite auf der \emph{südlichen} Hemisphäre
ähnliche Beobachtungen der Sonne und südlicher Circumpolarsterne mit
\textsc{Reichenbach}schen Repetitionskreisen von einem geübten Beobachter
angestellt würden. Für \emph{Reisen} in ferne Weltgegenden, zu astronomischen
Zwecken, wodurch früher so manche wichtige Erfahrungen und Aufschlüsse uns
geworden sind, ist in der That von den Regierungen seit geraumer Zeit gar
nichts geschehen, so sehr dieselben sonst auch in Europa die Wissenschaft
unterstützen.
